% =============================================================================
% Capítulo 8 — Conclusão
% =============================================================================
\chapter{Conclusão}
\label{cap:conclusao}

A criação de comportamentos de NPCs em jogos digitais enfrenta uma tensão
recorrente entre previsibilidade e adaptabilidade. Máquinas de Estados Finitos
oferecem controle total e facilidade de depuração, mas produzem oponentes cujos
padrões se tornam evidentes após poucas interações. O Aprendizado por Reforço
permite a emergência de comportamentos adaptativos, mas impõe o desafio de
definir funções de recompensa que orientem o agente sem recorrer a modelagem
manual exaustiva ou à imitação direta de demonstrações humanas, que vincula o
agente ao comportamento observado sem preservar a rastreabilidade causal das
decisões de \textit{design}.

O presente trabalho endereçou essa tensão por meio de uma abordagem que
integra Proveniência de Dados e Aprendizado por Reforço. Foi desenvolvido o
\textit{vertical slice} do jogo \textit{Chronicles Of The Lost Word}, um
protótipo 2D em Unity centrado em uma arena de combate contra um chefe. O
pipeline implementado percorreu todas as etapas propostas: sessões jogadas
contra o chefe FSM geraram grafos de proveniência com cadeias causais de
eventos; um processo automatizado extraiu sequências de ações associadas a
desfechos vitoriosos; essas sequências foram carregadas por um módulo de
\textit{Reward Shaping} que concedeu recompensa adicional ao agente PPO
durante o treinamento via Unity ML-Agents; e o modelo resultante foi exportado
em formato ONNX e integrado ao jogo em modo de inferência para avaliação
quantitativa.

Os resultados da avaliação --- 10 sessões contra o chefe FSM e 10 contra o
chefe PPO, com métricas coletadas automaticamente --- confirmaram que a
metodologia produziu um agente com comportamento mensuravelmente distinto do
\textit{baseline}. O chefe PPO gerou sessões de combate 2,5 vezes mais longas,
executou 53 vezes mais ações de ataque por sessão e foi a única condição
experimental a produzir uma derrota do jogador na amostra. As ações
participantes do mecanismo de \textit{Reward Shaping} --- ataque, pulo e
\textit{dash} --- corresponderam a 77,48\% das decisões do agente, indicando
que o sinal auxiliar derivado da proveniência exerceu influência direta sobre
a política aprendida.

Esses resultados validam a tese central do trabalho: sequências extraídas de
grafos de proveniência podem ser utilizadas como sinal auxiliar de recompensa
em um agente PPO, orientando o aprendizado sem imitação direta e produzindo
um comportamento que, embora distinto do humano, é rastreável às decisões
causais que o originaram. O agente não reproduziu trajetórias gravadas nem
utilizou discriminadores adversários; ele explorou o ambiente livremente e
convergiu para uma política influenciada por padrões causais vitoriosos ---
uma via alternativa ao GAIL \cite{ho2016gail} que preserva a interpretabilidade
do processo de treinamento.

As limitações observadas --- em particular, a repetição mecânica de ações
decorrente da ausência de \textit{cooldowns} no espaço de ações do agente e
do modelo congelado sem retreino após ajustes na avaliação --- não invalidam
a contribuição; elas a contextualizam. O comportamento repetitivo é uma
consequência documentada e explicável de escolhas de projeto específicas, e
sua correção constitui um passo natural de refinamento, não uma refutação da
abordagem. A amostra reduzida e a avaliação conduzida pelo próprio
desenvolvedor delimitam o alcance das conclusões à demonstração de viabilidade,
sem pretensão de generalização estatística.

Em última análise, o trabalho demonstrou que a proveniência de dados ---
originalmente concebida para rastreamento de linhagem e análise
\textit{post-mortem} de sessões --- pode ser reaproveitada como ponte entre
o \textit{gameplay} humano e o treinamento de agentes inteligentes em jogos.
Essa integração abre um espaço de investigação no qual grafos causais de
sessões de jogo deixam de ser apenas ferramentas de análise e passam a atuar
como fontes ativas de orientação para o aprendizado, combinando a
rastreabilidade da proveniência com a capacidade adaptativa do aprendizado
por reforço.
